\documentclass[11pt]{article}
\usepackage[margin=1in]{geometry}
\usepackage{hyperref}
\usepackage{parskip}
\usepackage{xcolor}

\hypersetup{
    colorlinks=true,
    linkcolor=blue,
    urlcolor=blue,
    citecolor=blue
}

\begin{document}

\noindent 13 August 2026

\vspace{1em}

\noindent The Editors\\
\textbf{Journal of Computer Languages}\\
Elsevier

\vspace{1.5em}

\noindent Dear Editors,

I submit for your consideration my manuscript, \textbf{``Grammar-Derived Operation Abstraction for Cross-Project Vulnerability Detection''}, for publication as a research article in the \emph{Journal of Computer Languages}.

\vspace{1em}
\noindent\textbf{What the paper is about.}
The paper studies a program abstraction: a labelling function over the node alphabet of a C grammar, and the question of what a node in a program graph should carry when a learned analysis is expected to transfer to code it has never seen. I compare, under strict experimental control, three node representations --- project-specific lexical tokens, an engineered taxonomy of operation types constructed from API allow-lists and identifier rules, and the parser's own grammar node kinds.

\vspace{1em}
\noindent\textbf{Why I believe it fits this journal.}
The object of study is a language-level artefact rather than a machine-learning architecture. I specify the abstraction as a total function over a tree-sitter grammar alphabet, give the projection lattice relating its granularities with a nesting proof, and --- unusually for this application area --- evaluate it \emph{as a static analysis}, reporting per-class precision and recall against annotated ground truth. That evaluation is not decoration: it uncovered three defects invisible to downstream accuracy, including a rule that classified pointer writes as deallocations and a class that counted 655 stream-I/O call sites as memory operations.

\vspace{1em}
\noindent\textbf{The principal finding.}
The representation matters more than the architecture, and the representation that works is the one that costs nothing. Grammar node kinds improve cross-project ranking by $+0.078$ AUC ($p = 0.007$) and $+0.071$ AUPRC ($p = 0.002$) over lexical tokens, winning every project-level fold, whereas the three message-passing schemes I compare differ by less than half as much under a fixed representation. The engineered taxonomy --- the artefact I built first and expected to be the contribution --- is statistically indistinguishable from lexical features ($p \approx 0.9$). I report this negative result at the same length as the positive one, and explain it: an information-theoretic analysis shows the taxonomy is $98.2\%$ predictable from the grammar kind alone ($H(\phi \mid \kappa) = 0.06$ bits), because most of its rules are syntax-directed and therefore already decided by the parser. I recommend that redundancy check as routine practice before an abstraction is credited with an effect.

\vspace{1em}
\noindent\textbf{Methodological care.}
Cross-project evaluation in this area is easy to get wrong, so I control for the failure modes documented in the literature: near-duplicate functions are removed before splitting; projects are the unit of partition, with repositories too large to place without dominating a fold held train-only and reported; vocabularies are fitted on training projects only; and confidence intervals come from a cluster bootstrap that resamples held-out projects rather than functions. Every table reports the trivial all-positive classifier as a floor.

\vspace{1em}
\noindent\textbf{Originality and availability.}
The manuscript is original, is not under consideration elsewhere, and has not been published previously. All code, annotations, and the scripts that regenerate every reported number are publicly available.

I confirm that the manuscript has been approved for submission, that it is not under consideration elsewhere, and that I have no competing interests to declare.

Thank you for your consideration.

\vspace{2em}

\noindent Yours sincerely,

\vspace{1.5em}

\noindent\textbf{Quang-Vinh Dang}\\
British University Vietnam, Hung Yen, Vietnam\\
Email: \texttt{vinh.dq4@buv.edu.vn}\\
ORCID: \href{https://orcid.org/0000-0002-3877-8024}{0000-0002-3877-8024}

\end{document}
